\subsection{مبانی نظری و فرمول‌بندی ریاضی}
فرض کنید لایه خطی مورد نظر دارای ماتریس وزن $\mathbf{W} \in \mathbb{R}^{d_{\text{out}} \times d_{\text{in}}}$ با تعداد پارامتر اولیه $P_{\text{orig}} = d_{\text{in}} \cdot d_{\text{out}}$ باشد.

\begin{enumerate}
    \item \textbf{تخصیص بودجه و تعیین رتبه مدوله‌شده:}
    اگر نسبت بودجه مجاز برای پارامترهای لایه پس از فشرده‌سازی با $\beta \in (0, 1)$ (\lr{total\_param\_ratio}) مشخص گردد، بودجه مجاز پارامترها برابر است با:
    \begin{equation}
    P_{\text{allowed}} = \beta \cdot (d_{\text{in}} \cdot d_{\text{out}})
    \end{equation}
    در تجزیه به ساختار دو فاکتوره با رتبه $r$، تعداد پارامترهای حاصله برابر است با $P_{\text{fact}} = r \cdot (d_{\text{in}} + d_{\text{out}})$. از این رو، رتبه بهینه مدوله‌شده به صورت مستقیم از رابطه زیر تعیین می‌شود:
    \begin{equation}
    r_{\text{target}} = \left\lfloor \frac{P_{\text{allowed}}}{d_{\text{in}} + d_{\text{out}}} \right\rfloor = \left\lfloor \frac{\beta \cdot d_{\text{in}} d_{\text{out}}}{d_{\text{in}} + d_{\text{out}}} \right\rfloor
    \end{equation}

    \item \textbf{اعمال محدودیت‌های مرزی و سقف کاهش پارامتر:}
    جهت تضمین کاهش قطعی حجم و ممانعت از تنزل بیش از حد رتبه، مقدار $r$ با روابط زیر مهار می‌گردد:
    \begin{equation}
    r_{\text{cap}} = \max\left(1, \left\lfloor \frac{d_{\text{in}} d_{\text{out}}}{d_{\text{in}} + d_{\text{out}}} \right\rfloor - 1\right)
    \end{equation}
    \begin{equation}
    r = \max\left(r_{\text{min}}, \min\left(r_{\text{target}}, \min(d_{\text{out}}, d_{\text{in}}) - 1, r_{\text{cap}}\right)\right)
    \end{equation}

    \item \textbf{تجزیه ماتریسی و ساخت لایه \lr{SVDLinear}:}
    با تثبیت رتبه مدوله‌شده $r$، تجزیه مقادیر منفرد بر روی ماتریس وزن اعمال شده و دو فاکتور $\mathbf{V}_{\text{down}} \in \mathbb{R}^{r \times d_{\text{in}}}$ و $\mathbf{U}_{\text{up}} \in \mathbb{R}^{d_{\text{out}} \times r}$ استخراج می‌شوند:
    \begin{equation}
    \mathbf{W} \approx \mathbf{U}_r \mathbf{\Sigma}_r \mathbf{V}_r^\top \implies \mathbf{V}_{\text{down}} = \mathbf{\Sigma}_r^{1/2} \mathbf{V}_r^\top, \quad \mathbf{U}_{\text{up}} = \mathbf{U}_r \mathbf{\Sigma}_r^{1/2}
    \end{equation}
    انتشار خروجی با رابطه خطی زنجیره‌ای بدون نیاز به ساخت ماتریس متراکم اولیه انجام می‌شود:
    \begin{equation}
    \mathbf{y} = \mathbf{x} \mathbf{V}_{\text{down}}^\top \mathbf{U}_{\text{up}}^\top + \mathbf{b}
    \end{equation}
\end{enumerate}


\paragraph{تحلیل طیفی و نامساوی‌های پایداری خطا}
برای بررسی میزان دقت تقریب ماتریس وزن یا نگاشت توجه در ماژول \lr{sdrm\_compressor}، تفاضل ماتریس فاکتورگیری‌شده از نگاشت متراکم اصلی بر اساس نامساوی هافدینگ-فون‌نیومان ارزیابی می‌گردد. فرض کنید $\mathbf{W} \in \mathbb{R}^{d_{\text{in}} \times d_{\text{out}}}$ ماتریس وزن هدف باشد که توسط تقریب کم‌مرتبه یا پراکنده با رتبه موثر $r$ جایگزین شده است. کران بالای خطای بازسازی بر مبنای نرم فروبنیوس به صورت زیر تصریح می‌شود:
\begin{equation}
\label{eq:spectral_bound_sdrm_compressor}
\| \mathbf{W} - \mathbf{W}_{r} \|_F = \sqrt{\sum_{i=r+1}^{\min(d_{\text{in}}, d_{\text{out}})} \sigma_i^2(\mathbf{W})} \le \epsilon \| \mathbf{W} \|_F
\end{equation}
که در آن $\sigma_i(\mathbf{W})$ بیانگر $i$-امین مقدار منفرد مرتب‌شده ماتریس به صورت نزولی است. نسبت حفظ انرژی طیفی $\eta_r = \frac{\sum_{i=1}^r \sigma_i^2}{\sum_{i=1}^{d} \sigma_i^2}$ تضمین می‌کند که با انتخاب آستانه $\tau \ge 0.90$، بیش از نود درصد از پراکندگی واریانس داده‌ها حفظ شده و افت سیگنال در انتشار گرادیان رو به عقب به حداقل ممکن کاهش می‌یابد.

\paragraph{تحلیل مجانبی و پیچیدگی اشتراک حافظه سخت‌افزاری}
از منظر پیچیدگی محاسباتی و انتقال داده‌ها در سلسومراتب حافظه (حافظه سراسری \lr{HBM} به حافظه محلی \lr{SRAM})، عملیات ضرب ماتریسی در این مکانیزم نیازمند تخصیص بافر موقت با ابعاد فاکتورگیری‌شده است. برای پردازش دنباله‌ای به طول $N$ با بعد مخفی $D$، بار محاسباتی هسته از فرم درجه دو به رابطه زیر تقلیل می‌یابد:
\begin{equation}
\label{eq:complexity_sdrm_compressor}
\mathcal{C}_{\text{sdrm\_compressor}} = \mathcal{O}\left( B \cdot N \cdot D \cdot r + B \cdot H \cdot N \cdot d_k \right) \ll \mathcal{O}\left( B \cdot H \cdot N^2 \cdot d_k \right)
\end{equation}
این بهبود مرتبه پیچیدگی، با کاهش چشمگیر حجم تبادل تانسورها در گذرگاه‌های ارتباطی کارت گرافیک، مانع از اشباع پهنای باند حافظه (پدیده \lr{Memory-Bound}) در فرآیند استنتاج طولانی‌مدت می‌گردد.
